\section{OBJETIVOS DO CURSO}
\label{OBJETIVOS}

\subsection{Objetivo Geral}
~
O curso tem como objetivo formar engenheiros capazes de atuar na concepção, análise, desenvolvimento, integração e gestão de sistemas de controle e automação, fundamentando-se no desenvolvimento de competências profissionais que articulem conhecimentos científicos e tecnológicos, habilidades de aplicação prática e atitudes orientadas à responsabilidade técnica e social. Pretende-se preparar um profissional apto a atuar em contextos industriais dinâmicos, a assumir funções de coordenação e gestão de processos e equipes, bem como a desenvolver atividades de pesquisa aplicada e docência, caso opte pela continuidade de sua formação acadêmica.

\subsection{Objetivos Específicos}
~
A formação no \NOMECURSO organiza-se segundo o desenvolvimento integrado de competências, compreendidas em três dimensões interdependentes:

\begin{itemize}
	\item \textbf{Conhecimentos}: fundamentos de matemática, física, computação, eletrônica, controle, automação, gestão e produção industrial, que sustentam a compreensão dos princípios que regem a concepção e a operação de sistemas automatizados;
	\item \textbf{Habilidades}: capacidade de projetar, implementar, testar, manter e otimizar sistemas de controle e automação; analisar dados e indicadores de desempenho; aplicar normas técnicas; utilizar ferramentas digitais avançadas; comunicar-se com clareza em ambientes técnicos e acadêmicos; coordenar equipes e processos;
	\item \textbf{Atitudes}: postura ética, responsabilidade socioambiental, visão crítica e sistêmica, iniciativa, criatividade, liderança colaborativa, capacidade de investigação científica e comprometimento com a atualização profissional contínua.
\end{itemize}

Com base nessas dimensões, os objetivos específicos do curso são:

\begin{itemize}
	\item Desenvolver competência científica e tecnológica para atuação qualificada no controle de processos e na automação industrial;
	\item Capacitar o estudante para atuar em equipes multidisciplinares e para exercer funções de gestão de projetos, processos e pessoas no contexto industrial;
	\item Formar profissionais aptos a avaliar e tomar decisões fundamentadas, considerando critérios técnicos, econômicos, ambientais e de segurança;
	\item Incentivar a investigação aplicada, o desenvolvimento de soluções inovadoras e a participação em projetos de pesquisa, extensão e inovação;
	\item Estimular a possibilidade de continuidade acadêmica em programas de pós-graduação, favorecendo o desenvolvimento de competências para atuação como pesquisador e docente;
	\item Favorecer a integração do estudante com o setor produtivo mediante estágios supervisionados, projetos integradores e cooperação com arranjos produtivos locais e regionais;
	\item Orientar para o cumprimento rigoroso das normas técnicas, regulamentações de segurança, saúde e meio ambiente;
	\item Promover postura empreendedora, liderança e capacidade de gestão orientada por responsabilidade social e sustentabilidade;
	\item Assegurar formação flexível, permitindo a incorporação de novas tecnologias emergentes no campo do controle e automação.
\end{itemize}
